\chapter{Wprowadzenie i analiza problemu}
\label{cha:wprowadzenie}

\section{Opis analizowanego problemu}
\label{sec:problem}

Sprawdzanie wiedzy uczniów za pomocą testów jest jednym z podstawowych zadań nauczyciela, jednak
w tradycyjnej, papierowej formie jest ono czasochłonne i mało wygodne. Wymaga ręcznego przygotowania
arkuszy, ich powielenia, a następnie żmudnego sprawdzania odpowiedzi i podliczania punktów. Dodatkowo
gromadzenie wyników oraz analiza tego, które zagadnienia sprawiają uczniom najwięcej trudności, są
w takim modelu praktycznie niemożliwe do przeprowadzenia na bieżąco.

Istnieją platformy do testów online, jednak najczęściej są to rozwiązania webowe, wymagające stałego
dostępu do Internetu i kont w usługach zewnętrznych, a ich pełne możliwości są zwykle płatne. Brakuje
prostego, samodzielnego narzędzia desktopowego, które łączyłoby tworzenie testów, ich automatyczne
ocenianie oraz analizę wyników w jednym miejscu, z danymi przechowywanymi w lokalnej bazie.

\section{Cel projektu}
\label{sec:cel}

Celem projektu było stworzenie od podstaw funkcjonalnej aplikacji okienkowej w języku C\#,
współpracującej z relacyjną bazą danych, która realizuje pełny cykl pracy z testem:

\begin{itemize}
  \item rejestrację i logowanie użytkowników z bezpiecznym przechowywaniem haseł oraz rozróżnieniem
        ról nauczyciela i ucznia,
  \item tworzenie testów złożonych z pytań jednokrotnego wyboru, wielokrotnego wyboru i pytań
        otwartych, z możliwością sterowania ich widocznością,
  \item przypisywanie testów grupom uczniów,
  \item rozwiązywanie testów przez uczniów w ograniczonym czasie, z obsługą wielu podejść,
  \item automatyczne ocenianie pytań zamkniętych oraz ręczne ocenianie i korektę punktów,
  \item wykrywanie prób oszustwa podczas rozwiązywania testu,
  \item analizę wyników w postaci statystyk oraz eksport do pliku CSV.
\end{itemize}

\section{Zakres realizowanych funkcjonalności}
\label{sec:zakres}

Zakres projektu obejmuje następujące moduły:

\begin{enumerate}
  \item \textbf{Moduł uwierzytelniania} - rejestracja (z wyborem roli), logowanie, haszowanie haseł,
        rozróżnienie ról nauczyciel/uczeń.
  \item \textbf{Moduł zarządzania grupami} - tworzenie grup, zarządzanie ich składem, przynależność
        ucznia do wielu grup.
  \item \textbf{Moduł testów} - tworzenie, edycja, duplikowanie i usuwanie testów, edytor pytań,
        ustawienia (czas, próg, liczba podejść), sterowanie widocznością (aktywny/nieaktywny).
  \item \textbf{Moduł rozwiązywania testów} - rozwiązywanie w ograniczonym czasie, nawigacja między
        pytaniami, limit podejść, wykrywanie opuszczenia okna testu.
  \item \textbf{Moduł oceniania} - automatyczne ocenianie pytań zamkniętych, ręczne ocenianie pytań
        otwartych oraz korekta oceny dowolnego podejścia.
  \item \textbf{Moduł statystyk} - zdawalność, średni wynik, trudność pytań, tabela wyników,
        eksport do CSV.
\end{enumerate}

\section{Przegląd istniejących rozwiązań}
\label{sec:przeglad}

Przed realizacją projektu przeanalizowano kilka istniejących narzędzi o podobnym przeznaczeniu.

\subsection{Kahoot!}

Popularna platforma do quizów i testów, działająca w przeglądarce, często wykorzystywana w szkołach.
Jej zaletą jest grywalizacja i prostota, jednak jest to wyłącznie usługa internetowa, wymagająca
stałego dostępu do sieci i konta, a obsługuje głównie pytania zamknięte. Pełne funkcje (m.in.
rozbudowane raporty) są dostępne w wersji płatnej.

\subsection{Microsoft Forms / Google Forms}

Webowe narzędzia do tworzenia formularzy i quizów, umożliwiające automatyczne sprawdzanie odpowiedzi
zamkniętych i zbieranie wyników. Są wygodne i darmowe w podstawowym zakresie, ale stanowią ogólne
narzędzia do ankiet - brakuje im funkcji typowo dydaktycznych, takich jak limit podejść, grupy
uczniów z przypisanymi testami, ręczne ocenianie pytań otwartych z korektą czy wykrywanie prób
oszustwa. Wymagają też konta w usłudze chmurowej i połączenia z Internetem.

\subsection{Moduł testów w systemie Moodle}

Moodle to rozbudowana platforma e-learningowa oferująca zaawansowany moduł testów (różne typy pytań,
banki pytań, statystyki). Jest jednak ciężka we wdrożeniu - wymaga serwera, bazy danych i konfiguracji
całej platformy - co czyni ją nieproporcjonalnie złożoną do prostego scenariusza, w którym nauczyciel
chce jedynie szybko ułożyć test i sprawdzić wyniki na własnym komputerze.

\subsection{Wnioski z analizy}

Żadne z analizowanych rozwiązań nie łączy lekkiej, samodzielnej aplikacji desktopowej z własną,
lokalną bazą danych, automatycznym i ręcznym ocenianiem (w tym pytań otwartych), obsługą grup,
limitem podejść, sterowaniem widocznością testu oraz wykrywaniem prób oszustwa - bez konieczności
korzystania z usług chmurowych. Te cechy wyznaczyły zakres tego projektu.
